\section{Tipo de Estudo}

Trata-se de um projeto tecnológico aplicado: o desenvolvimento de uma aplicação para
adquirir, processar e interpretar dados de composição corporal obtidos por bioimpedância
segmentada. A proposta integra dispositivos de medição corporal, processamento de dados e
técnicas de inteligência artificial voltadas à composição corporal.

Além do desenvolvimento do sistema, foram realizados testes funcionais e de integração da
plataforma, descritos na Seção~\ref{sec:testes}. As avaliações com usuários (usabilidade)
e a validação preliminar das interpretações geradas pelo modelo de inteligência artificial
constituem o protocolo planejado e não foram executadas nesta etapa, sendo tratadas como
trabalho futuro (Capítulo~\ref{cap:trabalhosfuturos}).

\section{Equipamento e Princípio de Funcionamento}

\subsection{Equipamento Utilizado}

O equipamento central do projeto é a balança de bioimpedância CF577BLE, fabricada pela
Unique Body Weight High. É um analisador corporal pessoal baseado em Análise de Impedância
Bioelétrica (BIA), capaz de medir a composição corporal de forma segmentada. O equipamento
foi escolhido pela arquitetura de oito eletrodos e pela variedade de parâmetros corporais
disponíveis em sua interface de programação. Ao contrário das balanças domésticas comuns,
a CF577BLE fornece dados segmentados, com análises separadas de membros superiores,
membros inferiores e tronco. O aparelho também tem conectividade Bluetooth Low Energy
(BLE), que sincroniza os dados com plataformas externas, característica essencial para a
integração com o \sintonia.

\subsection{Princípio de Funcionamento da Balança CF577BLE}

O funcionamento da balança baseia-se nos princípios da BIA apresentados na
Seção~\ref{sec:principios}. Durante a medição, correntes elétricas de baixa intensidade
são aplicadas ao corpo por oito eletrodos distribuídos entre a plataforma de apoio dos pés
e um manípulo segurado pelo usuário. Com isso, o equipamento faz análises segmentadas: as
correntes percorrem membros inferiores, membros superiores e tronco, e o aparelho estima
separadamente os parâmetros de braços, pernas e tronco, reduzindo limitações comuns nos
sistemas de apenas quatro eletrodos.

A arquitetura multifrequencial da CF577BLE mede em diferentes frequências, o que dá mais
sensibilidade na avaliação da água intra e extracelular. A partir dos valores de
impedância, algoritmos proprietários do \emph{firmware} estimam os indicadores corporais,
entre eles peso, IMC, frequência cardíaca, taxa metabólica basal, percentual de gordura
total e segmentado, massa muscular total e segmentada, gordura visceral, massa óssea, água
corporal total, água intracelular e extracelular, massa proteica, idade corporal e índice
geral de saúde. Após a aquisição, as informações são transmitidas por BLE. Embora o
ecossistema original do fabricante (Lefu/Unique Health) sincronize esses dados com a nuvem,
o \sintonia\ opta por ler a balança diretamente por BLE, sem intermediários.

\subsection{Avaliação Experimental e Adequação da CF577BLE para a População Brasileira}

O trabalho previa uma etapa exploratória para avaliar o comportamento das estimativas da
CF577BLE em uma amostra preliminar de voluntários brasileiros, verificando a
compatibilidade dos resultados com modelos descritos na literatura nacional. Essa
avaliação empírica com voluntários não foi conduzida nesta etapa
(Capítulo~\ref{cap:trabalhosfuturos}); descrevem-se a seguir o protocolo planejado e os
modelos efetivamente aplicados no processamento.

A motivação decorre de evidências de vieses sistemáticos em dispositivos de bioimpedância
quando usados em populações diferentes das de calibração dos algoritmos proprietários,
o que afeta as estimativas de massa livre de gordura, gordura corporal e distribuição
hídrica. Os resultados de Wahrlich et al.~\cite{wahrlich2026} são uma das principais
referências do trabalho: foi empregada a equação multivariada que os autores propõem para
estimar a FFM, desenvolvida e validada na comparação entre bioimpedância e DXA. Como este
estudo não dispõe de medições próprias por DXA, a equação é usada como base para processar
e interpretar os dados da CF577BLE, em uma investigação exploratória de sua aplicabilidade
em uma plataforma voltada à população brasileira.

A análise comparativa dos dados obtidos frente aos modelos da literatura, para identificar
diferenças sistemáticas ligadas à população de calibração, às frequências e aos modelos
proprietários, dependia da coleta com voluntários e não foi realizada nesta etapa,
permanecendo como trabalho futuro (Capítulo~\ref{cap:trabalhosfuturos}). A investigação de
ajustes matemáticos adicionais (regressão linear e modelos multivariados) para reduzir
diferenças específicas entre a CF577BLE e modelos de referência também não foi conduzida,
ficando registrada como trabalho futuro. Nesta versão, o sistema aplica diretamente a
equação de Wahrlich et al.~\cite{wahrlich2026}, desenvolvida e validada para a população
brasileira. Por se tratar de abordagem exploratória e de amostra reduzida, eventuais
resultados dessa natureza não teriam finalidade diagnóstica nem caráter de validação
clínica definitiva.

\subsubsection{Avaliação Comparativa com Métodos de Referência}

Considerando as limitações da bioimpedância, estava prevista uma análise comparativa dos
resultados da CF577BLE com valores de referência da literatura. O objetivo não seria
validar clinicamente o equipamento, mas verificar a coerência fisiológica das estimativas
usadas pelo \sintonia. Seriam observados indicadores como massa livre de gordura,
percentual de gordura corporal, água corporal total e taxa metabólica basal, confrontados
com modelos publicados. Essa análise permitiria identificar desvios sistemáticos e
subsidiar futuras estratégias de calibração computacional (Capítulo~\ref{cap:trabalhosfuturos}).

\subsection{Avaliação das Frequências Elétricas e Compatibilidade com Modelos Clínicos}

Equipamentos clínicos muito usados em pesquisa, como a Tanita BC-418, medem em frequências
próximas de 50~kHz. A CF577BLE emprega arquitetura multifrequencial com medições em
frequências distintas, incluindo 20~kHz e 100~kHz. Essa diferença pode influenciar as
estimativas, pois frequências distintas têm sensibilidades diferentes para os
compartimentos de água intra e extracelular; assim, modelos desenvolvidos para uma
frequência podem não se comportar igual em dispositivos multifrequenciais.

Quanto à compatibilização entre frequências, foi implementada a interpolação logarítmica
da impedância medida em 20~kHz e 100~kHz para uma impedância equivalente em 50~kHz
(Seção~\ref{sec:equacoes}, Equação~\ref{eq:interp50}), o que permite aplicar modelos
desenvolvidos para 50~kHz. Estratégias adicionais de compatibilização por regressão e
adaptações inspiradas em arquiteturas de equipamentos comerciais não foram desenvolvidas
nesta etapa, permanecendo como trabalho futuro (Capítulo~\ref{cap:trabalhosfuturos}). Essas
estratégias adicionais de compatibilização entre frequências e de calibração computacional
são retomadas no Capítulo~\ref{cap:trabalhosfuturos}.

\section{Coleta de Dados}

A coleta de dados com usuários voluntários não foi realizada nesta etapa do trabalho; o
protocolo planejado está reposicionado como trabalho futuro
(Capítulo~\ref{cap:trabalhosfuturos}). As subseções a seguir resumem o que estava previsto
e descrevem o tratamento e os aspectos éticos dos dados.

\subsection{Participantes}

A etapa de coleta de dados com usuários voluntários, prevista para validar a solução na
plataforma \sintonia, não foi conduzida como estudo de campo nesta etapa. O planejamento
previa uma amostra de conveniência de aproximadamente 15 a 30 usuários, e foi
reposicionada como trabalho futuro (Capítulo~\ref{cap:trabalhosfuturos}).

\subsection{Procedimento de Coleta}

O procedimento de coleta previsto para esse estudo de campo (cadastro dos participantes,
leitura via Web Bluetooth e registro dos parâmetros) está descrito no
Capítulo~\ref{cap:trabalhosfuturos}.

\subsection{Armazenamento e Tratamento dos Dados}

Os dados coletados seriam armazenados em banco estruturado, permitindo construir o
histórico individual de cada participante. Antes do uso para análises estatísticas ou
treinamento dos modelos, os dados seriam submetidos a validação, padronização e
anonimização, removendo informações que permitissem a identificação direta. As informações
poderiam ser usadas para análise temporal da composição corporal, avaliação de desempenho
e aperfeiçoamento dos modelos de IA, respeitando princípios de anonimização e uso
responsável~\cite{etic2024}.

\subsection{Aspectos Éticos e Privacidade}

Como o projeto envolve dados de saúde e composição corporal, seriam adotadas medidas de
proteção da privacidade e de segurança das informações, seguindo recomendações recentes
para aplicações de IA em saúde~\cite{etic2024}. Os dados de pesquisa seriam tratados de
forma anonimizada, com acesso restrito aos pesquisadores, e o armazenamento seguiria
práticas compatíveis com saúde digital. O \sintonia\ não tem finalidade diagnóstica e não
substitui avaliação de profissionais de saúde: suas interpretações têm caráter informativo
e educacional.

\section{Avaliação da Solução Proposta}

A avaliação do \sintonia\ foi planejada em três frentes: usabilidade, desempenho técnico e
qualidade das interpretações do modelo de IA. Nesta etapa foram executados os testes
funcionais e de desempenho técnico relatados na Seção~\ref{sec:testes}. As avaliações com
usuários (escala SUS), a validação por especialistas e a análise sistemática das
interpretações do SLM, descritas a seguir, correspondem ao protocolo planejado e foram
reposicionadas como trabalho futuro (Capítulo~\ref{cap:trabalhosfuturos}).

\subsection{Avaliação de Usabilidade}

Como instrumento de avaliação de usabilidade, foi selecionada a escala \emph{System
Usability Scale} (SUS), amplamente utilizada para avaliação de sistemas
computacionais~\cite{brooke1996}. O instrumento permite obter indicadores quantitativos
sobre a percepção dos usuários quanto à facilidade de uso, clareza da interface e
satisfação geral. Sua aplicação com usuários, porém, não foi realizada nesta etapa
(Capítulo~\ref{cap:trabalhosfuturos}). Caso essa etapa seja realizada, os resultados
poderiam identificar oportunidades de melhoria na interface e fornecer evidências sobre a
aceitação da solução. Por depender de uma amostra de usuários voluntários ainda não
recrutada, a aplicação do questionário SUS foi reposicionada como trabalho futuro e não
constitui requisito desta etapa de validação técnica.

\subsection{Avaliação Técnica do Sistema}

A avaliação técnica contemplou a análise do funcionamento dos principais componentes da
arquitetura proposta, incluindo a leitura da balança via Web Bluetooth, o armazenamento dos
dados, o processamento das métricas corporais e a geração das interpretações automatizadas
(Seção~\ref{sec:testes}). Foram observados aspectos como a estabilidade da comunicação, a
consistência dos dados recebidos, a integridade dos registros armazenados e o tempo de
processamento. A Figura~\ref{fig:fluxoslm} apresenta o fluxo do SLM no \sintonia.

\begin{figure}[!ht]
    \centering
    \begin{tikzpicture}[
        node distance=6mm,
        proc/.style={draw, rounded corners, align=center, text width=.80\linewidth,
                     inner sep=6pt, font=\small},
        loc/.style={proc, fill=green!8},
        seta/.style={-{Latex[length=2.5mm]}, thick},
    ]
        \node[proc] (in) {\textbf{Entrada} \\ indicadores biométricos calculados, contexto da anamnese e histórico do usuário};
        \node[proc, below=of in] (prompt) {\textbf{Montagem do \emph{prompt}} \\ \emph{system prompt} (persona ``Vida'') $+$ regras de interpretação fundamentadas na literatura $+$ dados da medição};
        \node[loc, below=of prompt] (inf) {\textbf{Inferência local} \\ SLM Qwen 2.5 (7B) via Ollama, em servidor com GPU; os dados de saúde não saem do servidor (LGPD)};
        \node[proc, below=of inf] (pos) {\textbf{Pós-processamento} \\ trava de idioma (português), remoção de emojis e de notação não renderizável e verificação de procedência dos dados};
        \node[proc, below=of pos] (out) {\textbf{Saída} \\ interpretação em linguagem natural, apresentada no aplicativo web ou no WhatsApp};
        \draw[seta] (in) -- (prompt);
        \draw[seta] (prompt) -- (inf);
        \draw[seta] (inf) -- (pos);
        \draw[seta] (pos) -- (out);
    \end{tikzpicture}
    \caption{Fluxo do \emph{Small Language Model} no \sintonia. A especialização é feita por
    engenharia de \emph{prompt} e regras de interpretação, com inferência local e sem
    \emph{fine-tuning} do modelo (o treinamento supervisionado é tratado como trabalho futuro,
    Capítulo~\ref{cap:trabalhosfuturos}).}
    \label{fig:fluxoslm}
\end{figure}

\subsection{Avaliação das Interpretações Geradas pelo SLM}

Para avaliar as interpretações produzidas pelo \emph{Small Language Model}, definiu-se uma
análise exploratória com base em conceitos e recomendações da literatura sobre
bioimpedância e composição corporal. Essa avaliação sistemática não foi conduzida nesta
etapa, sendo reposicionada como trabalho futuro (Capítulo~\ref{cap:trabalhosfuturos}); a
verificação realizada limitou-se aos testes funcionais descritos na Seção~\ref{sec:testes}.
A avaliação qualitativa poderá contemplar aspectos como coerência das respostas,
consistência das recomendações, adequação das interpretações aos parâmetros fornecidos e
capacidade de identificar padrões relevantes, conforme abordagens de estudos com modelos
especializados de linguagem~\cite{abdin2024,touvron2024}. Caso essa etapa seja realizada,
as análises verificariam a aderência das respostas às regras de interpretação e às
referências científicas adotadas.

\subsection{Validação por Especialistas}

A validação humana especializada é recomendada em aplicações de IA voltadas à saúde,
especialmente em sistemas que produzem interpretações direcionadas aos
usuários~\cite{etic2024}. O escopo desta etapa concentra-se na validação técnica da
plataforma, na implementação dos modelos de processamento biométrico e na verificação da
aderência das interpretações às referências adotadas. Como aprofundamento futuro, poderá
ser realizada a submissão das interpretações à avaliação de profissionais de saúde,
educação física e nutrição, com critérios como relevância clínica, clareza das explicações,
consistência científica e adequação das recomendações. Por não constituir objetivo
principal desta fase, a validação por especialistas é considerada evolução futura do
\sintonia.

\subsection{Indicadores de Desempenho}

Para complementar a avaliação, foram analisados indicadores quantitativos de desempenho do
sistema, incluindo tempo médio de resposta, disponibilidade da plataforma, taxa de
processamento bem-sucedido das medições e consistência das interpretações
(Seção~\ref{sec:testes}). Esses indicadores permitem verificar a viabilidade técnica da
arquitetura proposta e subsidiar futuras evoluções da plataforma.
